\documentclass[10pt,a4paper]{amsart}
\usepackage[margin=19mm]{geometry}
\usepackage{amsmath,amssymb,mathtools}
\usepackage[scr=rsfs,cal=euler]{mathalfa}
\usepackage{xfrac}
\usepackage[unicode,hidelinks,bookmarks=false]{hyperref}
\newcommand{\ie}{i.e.\ }
\newcommand{\cf}{cf.\ }
\newcommand{\resp}{respectively\ }
\newcommand{\Subsection}{\subsection}
\providecommand{\nobreakdash}{\nobreak\hbox{-}}
\setlength{\emergencystretch}{2em}
\multlinegap0pt
\usepackage{amssymb}
\usepackage{mathtools}
\usepackage{xcolor}
\newtheorem{lem}{Lemma}
\renewcommand{\ge}{\geqslant}
\renewcommand{\le}{\leqslant}
\title{Sharp Upper Bounds for Moments of Dedekind Zeta Functions}
\author{Benjamin Durkan, Nilmoni Karak, Kamalakshya Mahatab}
\date{Source: arXiv:2606.27516v3 (CC BY 4.0)}
\begin{document}
\maketitle

\noindent\emph{Source and licence.} Sharp Upper Bounds for Moments of Dedekind Zeta Functions. Version arXiv:2606.27516v3 (CC BY 4.0), distributed by arXiv under the Creative Commons Attribution 4.0 International licence, http://creativecommons.org/licenses/by/4.0/, as recorded in the arXiv licence metadata for this exact version and shown on the abstract page https://arxiv.org/abs/2606.27516v3. Full licence notice: \texttt{licenses/arxiv-2606.27516-CC-BY-4.0.txt}.\par\smallskip

\noindent\textit{Editorial scope.} Counted unit: the article's lem shifted-moment-trivial-bound (source0.tex lines 1156--1161). The article's complete original proof of it is delivered with it (source0.tex lines 1162--1280, one \texttt{proof} environment of 119 lines). No mathematics is written, repaired or re-derived: the statement and the proof are the article's own lines, in the article's own notation, and no retained line is substituted or re-flowed.  Retained uncounted as the source-local context the proof needs: lines 644--671; lines 1076--1084.  Nothing was omitted from any retained span, so the package carries no omission record.  No retained line contains a citation, so the wrapper carries no bibliography and no bibliography entry.  No retained line contains a figure environment, an image-inclusion command or drawing code, so no image is delivered and the package declares no asset dependency.  Editorial formatting only, in addition: an AMS \texttt{amsart} preamble that restates the article's own declarations and macros used by the retained lines, namely the environments \texttt{lem} on separate counters, together with the standard packages those lines need.  The retained environments print in this document as Lemma 1 in order of appearance, whereas the article prints the same items with the numbering of its own full document.  The complete native source of the article as distributed in the arXiv e-print for this exact version is bundled unchanged as \texttt{sources/203771/source0.tex}.
\begin{lem}\label{complex-dirichlet-mean}
Let $\mathcal P_1,\ldots,\mathcal P_R$ be disjoint finite sets of primes and
let $\mathcal Q_1,\ldots,\mathcal Q_S$ be finite sets whose elements are
pairwise coprime, both within and between the sets, and coprime to every prime
in the $\mathcal P_r$.  Let
\begin{equation*}
D_r(t)=\sum_{p\in\mathcal P_r}c_r(p)p^{-it},\qquad
E_s(t)=\sum_{q\in\mathcal Q_s}d_s(q)q^{-it}.
\end{equation*}
Let $N_r,M_s$ be non-negative integers.  Suppose every integer occurring in
\begin{equation*}
 A(t):=\prod_{s=1}^S E_s(t)^{M_s}
 \prod_{r=1}^R\sum_{0\le m\le N_r}\frac{D_r(t)^m}{m!}
\end{equation*}
is at most $T^{1/10}$.  Then    
        \begin{align*}
         &\int_T^{2T}
        \left|\prod_{r=1}^{R}\sum_{0\le  m\le  N_r}\frac{D_r(t)^m}{m!}\right|^2
        \left|\prod_{s=1}^S E_s(t)^{M_s} \right|^2\,dt \\
&\ll T\exp\left(\sum_{r=1}^R\sum_{p\in\mathcal P_r}|c_r(p)|^2\right)
 \prod_{s=1}^S M_s!\left(\sum_{q\in\mathcal Q_s}|d_s(q)|^2\right)^{M_s}\\
&\quad+T^{1/10}
 \prod_{s=1}^S\left(\sum_{q\in\mathcal Q_s}|d_s(q)|\right)^{2M_s}
 \left(\prod_{r=1}^R\sum_{0\le m\le N_r}\frac1{m!}
 \left(\sum_{p\in\mathcal P_r}|c_r(p)|\right)^m\right)^2.
\end{align*}
Here empty products are interpreted as $1$.
\end{lem}

	\begin{lem}\label{conditional-upper-bound}
			Let $L/\mathbb Q$ be finite Galois and let $K$ be an intermediate field. Assume GRH for $\zeta_L(s)$. There is $T_0=T_0(K,L)\ge3$ such that, for $T\ge T_0$, uniformly for $T/2\le t\le5T/2$ and $e^2\le x\le T^2$, we have
		\begin{align*}
			&\log \left| \zeta_K\left(\tfrac{1}{2}+it\right)\right|
				\le  \Re\left( \sum_{p\le  x} \frac{\Lambda_K(p)}{p^{\frac{1}{2}+\frac{1}{\log x}+ it}\log p} \frac{\log\left(x/p\right)}{\log x}+ \sum_{p\le  \min\{ \sqrt{x}, \ \log T\}} \frac{\Lambda_K(p^2)}{2p^{1+2it}\log p}\right)\\
					& \hspace{3.5cm}+ \frac{[K:\mathbb{Q}]\log T}{\log x} +O_{K,L}(1),
		\end{align*}
		where $p$ denotes the rational prime.
	\end{lem}

\begin{lem}\label{shifted-moment-trivial-bound}
Let $K_1,\ldots,K_m$ be fixed number fields, let $L/\mathbb Q$ be a finite Galois extension containing them, and let $A_1,\ldots,A_m>0$ be fixed.  Assume GRH for $\zeta_L(s)$.  Uniformly for $|b_i|\le T/2$, there is a constant $C$, depending only on $L$, the fields, and the exponents, such that
		\begin{align*}
			\int_{T}^{2T} \prod_{j=1}^{m}\left| \zeta_{K_j}\left(\frac{1}{2}+i(t+b_j)\right)\right|^{A_j}\,dt\ll T(\log T)^C.
		\end{align*}
	\end{lem}

	\begin{proof}
It suffices to prove, for each fixed intermediate field $K$,  $A>0$ and $U\ge 3$ that
\begin{equation}\label{eq:one-field-crude-moment}
\int_U^{2U}\left|\zeta_K\left(\frac12+iu\right)\right|^A\,du
        \ll_{L,K,A} U(\log U)^{C(L,K,A)}.
\end{equation}
Indeed, GRH for $\zeta_L$ implies GRH for $\zeta_K$ by the Aramata--Brauer theorem.  Let $d=[K:\mathbb Q]$ and, for $V\ge3$, define
\begin{equation*}
        \mathcal E(V)=
        \left\{u\in[U,2U]:
        \log\left|\zeta_K\left(\frac12+iu\right)\right|>V\right\}.
\end{equation*}
We first establish the two estimates needed below.  There exist positive constants $C_1,C_2,C_3,D>0$, depending only on $K$ and $A$, such that, with
\begin{equation*}
        V_0=D\frac{\log U}{\log\log U},
\end{equation*}
one has
\begin{align}
\label{eq:dedekind-large-values-small}
        \operatorname{meas}(\mathcal E(V))
        &\ll
        U\exp\left(-\frac{V^2}{C_1\log\log U}\right)
         &\text{for }3\le  V\le  C_2\log\log U,\\
\label{eq:dedekind-large-values-large}
        \operatorname{meas}(\mathcal E(V))
        &\ll
        U\exp\left(-C_3V\log\frac{V}{\log\log U}\right)
        &\text{for }C_2\log\log U\le V\le V_0,
\end{align}
while $\mathcal E(V)=\emptyset$ for $V>V_0$.  To prove these assertions, let $x=e^Y$, where $2\le Y\le2\log U$.  By Lemma~\ref{conditional-upper-bound} and the bound $|\Lambda_K(n)|\le d\Lambda(n)$,
\begin{equation*}
        \sum_{p\le  \log U}\frac{|\Lambda_K(p^2)|}{p\log p}
        \ll_K \sum_{p\le\log U}\frac1p
        \ll\log\log\log U
\end{equation*}
gives, uniformly for $u\in[U,2U]$,
\begin{equation}\label{eq:sound-prime-majorant}
        \log\left|\zeta_K\left(\frac12+iu\right)\right|
        \le
        \Re P_x(u)+\frac{d\log U}{Y}+O_K(\log\log\log U),
\end{equation}
where
\begin{equation*}
        P_x(u)=
        \sum_{p\le  x}
        \frac{\Lambda_K(p)}{p^{1/2+1/Y+iu}\log p}
        \frac{\log(x/p)}{Y}.
\end{equation*}
If $r\ge1$ and $x^r\le U^{1/10}$, Lemma~\ref{complex-dirichlet-mean} gives
\begin{equation}\label{eq:Px-moment}
        \int_U^{2U}|P_x(u)|^{2r}\,du
        \ll
        U r!\left(C_K\log\log x\right)^r
        +U^{1/10}(C_Kx^{1/2})^{2r}
        \ll
        U r!\left(C_K\log\log x\right)^r+U^{1/5}C_K^{2r}.
\end{equation}
After division by $(V/2)^{2r}$, the off-diagonal term is at most
\begin{equation*}
        U^{1/5}\left(\frac{2C_K}{V}\right)^{2r}.
\end{equation*}
We now choose $Y=8d\log U/V$.  When $C_K\log\log\log U\le V\le C_2\log\log U$, membership in $\mathcal E(V)$ implies $\Re P_x(u)>V/2$.  With
\begin{equation*}
        r=\left\lfloor\frac{V^2}{C_4\log\log U}\right\rfloor+1,
\end{equation*}
and $C_4$ sufficiently large, we have $x^r\le U^{1/10}$, and Chebyshev's inequality gives
\begin{equation*}
        \operatorname{meas}(\mathcal E(V))
        \ll
        U\left(\frac{C_Kr\log\log U}{V^2}\right)^r
        \ll
        U\exp\left(-\frac{V^2}{C_1\log\log U}\right).
\end{equation*}
The range $3\le V<C_K\log\log\log U$ follows from $\operatorname{meas}(\mathcal E(V))\le U$ after enlarging $C_1$.  This proves \eqref{eq:dedekind-large-values-small}.  Lemma~\ref{conditional-upper-bound} with $x=(\log U)^2$ also gives
\begin{align*}
\log\left|\zeta_K\left(\frac12+iu\right)\right|
&\le d\sum_{p\le(\log U)^2}\frac1{\sqrt p}
+\frac{d\log U}{2\log\log U}+O_K(\log\log\log U)\\
&\le D\frac{\log U}{\log\log U}
\end{align*}
for a fixed $D=D(K)$; hence $\mathcal E(V)=\emptyset$ for $V>V_0$.  If $C_2\log\log U<V\le V_0$, keep the previous choice of $Y$ and take
\begin{equation*}
        r=\left\lfloor\frac{V}{C_5d}\right\rfloor,
\end{equation*}
where $C_5\ge80$ is fixed.  Then $Y\ge(8d/D)\log\log U$ and
\begin{equation*}
        rY\le\frac{8}{C_5}\log U,
\end{equation*}
so $x^r\le U^{1/10}$.  The diagonal term in \eqref{eq:Px-moment} gives
\begin{equation*}
        \operatorname{meas}(\mathcal E(V))
        \ll
        U\left(\frac{C_Kr\log\log U}{V^2}\right)^r
        \ll
        U\exp\left(-C_3V\log\frac{V}{\log\log U}\right),
\end{equation*}
after $C_2$ is enlarged.  In both ranges the off-diagonal contribution is $O(U^{1/5})$.  Taking $C_3D\le1/2$ makes the right side of \eqref{eq:dedekind-large-values-large} at least $U^{1/2}$, so this error is absorbed and \eqref{eq:dedekind-large-values-large} follows.

We now integrate the distribution estimates.  With $W(u)=\log|\zeta_K(1/2+iu)|$,
\begin{equation*}
        \int_U^{2U}e^{AW(u)}\,du
        \le
        Ue^{AC_2\log\log U}
        +A\int_{C_2\log\log U}^{\infty}
        e^{AV}\operatorname{meas}(\mathcal E(V))\,dV.
\end{equation*}
Choose $C_2$ so that $C_3\log C_2\ge A+2$.  The integral is then $O(U)$, proving \eqref{eq:one-field-crude-moment}.  Finally, let $A=A_1+\cdots+A_m$.  Hölder's inequality gives
\begin{equation*}
\int_T^{2T}\prod_{j=1}^m
\left|\zeta_{K_j}\left(\frac12+i(t+b_j)\right)\right|^{A_j}\,dt
\le
\prod_{j=1}^m
\left(
\int_{T+b_j}^{2T+b_j}
\left|\zeta_{K_j}\left(\frac12+iu\right)\right|^{A}\,du
\right)^{A_j/A}.
\end{equation*}
Each shifted interval $[T+b_j,2T+b_j]$ can be covered by $O(1)$ dyadic intervals whose endpoints are comparable to $T$.  Applying \eqref{eq:one-field-crude-moment} on each such interval completes the proof.
\end{proof}

\end{document}
